% GENERATED FILE. Edit canonical lessons, then run npm run generate:books.
% canonical-source-hash: fnv1a64:2347b0df32dacff6
% canonical-lessons: PA-C277-gahak, PA-C277-chunna, PA-C277-parakhna, PA-C277-chiz, PA-C277-saman

\chapter{Choosing the Best}
\label{ch:pa-a2-choosing-the-best}
\hlchaptermodality{277}

\begin{chapteropening}
\textbf{By the end of this chapter:} \emph{I can choose and check goods in a shop.}

Complete the last lesson of chapter 277: I can choose and check goods in a shop.
\end{chapteropening}

\section[\texorpdfstring{\pa{ਗਾਹਕ} (gāhak)}{ਗਾਹਕ (gāhak)}]{\pa{ਗਾਹਕ} (gāhak) — a customer}

\label{lesson:PA-C277-gahak}



\begin{quote}
Before the new one: say the Punjabi for a purchase, then the Punjabi for free of charge.
\end{quote}

\subsection*{You'll want to know: \pa{ਗਾਹਕ}}
\textbf{\pa{ਗਾਹਕ}} — \emph{gāhak} — \textquotedblleft{}a customer\textquotedblright{}.

\begin{quote}
\pa{ਸਵੇਰੇ} \pa{ਦੁਕਾਨ} \pa{ਵਿੱਚ} \pa{ਬਹੁਤ} \pa{ਗਾਹਕ} \pa{ਹੁੰਦੇ} \pa{ਹਨ।}
\end{quote}

\begin{quote}
\emph{savere dukān vich bahut gāhak hunde han.}
\end{quote}

\begin{quote}
In the morning there are many customers in the shop.
\end{quote}

A masculine or feminine noun, by the person meant.

\begin{cousinweb}
\textbf{\pa{ਗਾਹਕ}} (\emph{gāhak}) goes back to Sanskrit \emph{grāhaka}, \textquotedblleft{}one who takes; a purchaser\textquotedblright{} (Monier-Williams); CDIAL has no entry for it.
\end{cousinweb}

\begin{grammarlens}[title={What you've built}]
One more word for this chapter.
\end{grammarlens}

\subsection*{Your turn}
\begin{itemize}
\raggedright
  \item \emph{Say it:} \emph{gāhak}
  \item \emph{Say it:} \emph{gāhak}, once more
  \item \emph{Say it:} \emph{savere dukān vich bahut gāhak hunde han.} — In the morning there are many customers in the shop.
  \item \emph{Recall:} say the Punjabi for a purchase, then the Punjabi for free of charge, then say \emph{gāhak} again
\end{itemize}

\begin{tcolorbox}[breakable,colback=teal!4,colframe=teal!35!black,title={Before you move on}]
What does \pa{ਗਾਹਕ} mean? (\textquotedblleft{}A customer\textquotedblright{}.) Say it once more.
\end{tcolorbox}
\section[\texorpdfstring{\pa{ਚੁਣਨਾ} (chuṇnā)}{ਚੁਣਨਾ (chuṇnā)}]{\pa{ਚੁਣਨਾ} (chuṇnā) — to choose; to pick out}

\label{lesson:PA-C277-chunna}



\begin{quote}
Before the new one: say the Punjabi for free of charge, then the Punjabi for a customer.
\end{quote}

\subsection*{You'll want to know: \pa{ਚੁਣਨਾ}}
\textbf{\pa{ਚੁਣਨਾ}} — \emph{chuṇnā} — \textquotedblleft{}to choose; to pick out\textquotedblright{}.

\begin{quote}
\pa{ਮੈਂ} \pa{ਤਾਜ਼ੇ} \pa{ਸੇਬ} \pa{ਚੁਣਦਾ} \pa{ਹਾਂ।}
\end{quote}

\begin{quote}
\emph{maiṁ tāze seb chuṇdā hāṁ.}
\end{quote}

\begin{quote}
I pick out the fresh apples.
\end{quote}

Past (I …): \pa{ਮੈਂ} \pa{ਲਾਲ} \pa{ਕਮੀਜ਼} \pa{ਚੁਣੀ}, \emph{maiṁ lāl kamīz chuṇī}.

\begin{cousinweb}
\textbf{\pa{ਚੁਣਨਾ}} (\emph{chuṇnā}) is inherited from Sanskrit \emph{cinoti}, \textquotedblleft{}heaps up, gathers\textquotedblright{}; CDIAL lists Punjabi \emph{cuṇnā}, \textquotedblleft{}to build up in layers, plait\textquotedblright{}, under it, beside Lahnda \emph{cuṇaṇ}, \textquotedblleft{}to pick, gather, choose\textquotedblright{} (CDIAL 4814).
\end{cousinweb}

\begin{grammarlens}[title={What you've built}]
One more word for this chapter.
\end{grammarlens}

\subsection*{Your turn}
\begin{itemize}
\raggedright
  \item \emph{Say it:} \emph{chuṇnā}
  \item \emph{Say it:} \emph{chuṇnā}, once more
  \item \emph{Say it:} \emph{maiṁ tāze seb chuṇdā hāṁ.} — I pick out the fresh apples.
  \item \emph{Recall:} say the Punjabi for free of charge, then the Punjabi for a customer, then say \emph{chuṇnā} again
  \item \emph{Write it:} copy \textbf{\pa{ਚੁਣਨਾ}} once, checking the whole word
\end{itemize}

\begin{tcolorbox}[breakable,colback=teal!4,colframe=teal!35!black,title={Before you move on}]
What does \pa{ਚੁਣਨਾ} mean? (\textquotedblleft{}To choose; to pick out\textquotedblright{}.) Say it once more.
\end{tcolorbox}
\section[\texorpdfstring{\pa{ਪਰਖਣਾ} (parakhṇā)}{ਪਰਖਣਾ (parakhṇā)}]{\pa{ਪਰਖਣਾ} (parakhṇā) — to test; to check}

\label{lesson:PA-C277-parakhna}



\begin{quote}
Before the new one: say the Punjabi for a customer, then the Punjabi for to choose.
\end{quote}

\subsection*{You'll want to know: \pa{ਪਰਖਣਾ}}
\textbf{\pa{ਪਰਖਣਾ}} — \emph{parakhṇā} — \textquotedblleft{}to test; to check\textquotedblright{}.

\begin{quote}
\pa{ਮੇਰੀ} \pa{ਮਾਤਾ} \pa{ਹਰ} \pa{ਸੇਬ} \pa{ਪਰਖਦੀ} \pa{ਹੈ।}
\end{quote}

\begin{quote}
\emph{merī mātā har seb parakhdī hai.}
\end{quote}

\begin{quote}
My mother checks every apple.
\end{quote}

Past (I …): \pa{ਮੈਂ} \pa{ਨਵਾਂ} \pa{ਫ਼ੋਨ} \pa{ਪਰਖਿਆ}, \emph{maiṁ navāṁ fon parkhiā}.

\begin{cousinweb}
\textbf{\pa{ਪਰਖਣਾ}} (\emph{parakhṇā}) is inherited from Sanskrit \emph{parīkṣate}, \textquotedblleft{}observes, examines\textquotedblright{}; CDIAL lists Punjabi \emph{parakhṇā}, \textquotedblleft{}to test\textquotedblright{}, under it (CDIAL 7904).
\end{cousinweb}

\begin{grammarlens}[title={What you've built}]
One more word for this chapter.
\end{grammarlens}

\subsection*{Your turn}
\begin{itemize}
\raggedright
  \item \emph{Say it:} \emph{parakhṇā}
  \item \emph{Say it:} \emph{parakhṇā}, once more
  \item \emph{Say it:} \emph{merī mātā har seb parakhdī hai.} — My mother checks every apple.
  \item \emph{Recall:} say the Punjabi for a customer, then the Punjabi for to choose, then say \emph{parakhṇā} again
  \item \emph{Write it:} copy \textbf{\pa{ਪਰਖਣਾ}} once, checking the whole word
\end{itemize}

\begin{tcolorbox}[breakable,colback=teal!4,colframe=teal!35!black,title={Before you move on}]
What does \pa{ਪਰਖਣਾ} mean? (\textquotedblleft{}To test; to check\textquotedblright{}.) Say it once more.
\end{tcolorbox}
\section[\texorpdfstring{\pa{ਚੀਜ਼} (chīz)}{ਚੀਜ਼ (chīz)}]{\pa{ਚੀਜ਼} (chīz) — a thing; an item}

\label{lesson:PA-C277-chiz}



\begin{quote}
Before the new one: say the Punjabi for to choose, then the Punjabi for to test.
\end{quote}

\subsection*{You'll want to know: \pa{ਚੀਜ਼}}
\textbf{\pa{ਚੀਜ਼}} — \emph{chīz} — \textquotedblleft{}a thing; an item\textquotedblright{}.

\begin{quote}
\pa{ਇਹ} \pa{ਚੀਜ਼} \pa{ਬਹੁਤ} \pa{ਮਹਿੰਗੀ} \pa{ਹੈ।}
\end{quote}

\begin{quote}
\emph{eh chīz bahut mahiṅgī hai.}
\end{quote}

\begin{quote}
This item is very expensive.
\end{quote}

A feminine noun.

\begin{cousinweb}
\textbf{\pa{ਚੀਜ਼}} (\emph{chīz}) is a loan, usually traced to Persian \emph{chīz}, \textquotedblleft{}a thing\textquotedblright{}.
\end{cousinweb}

\begin{grammarlens}[title={What you've built}]
One more word for this chapter.
\end{grammarlens}

\subsection*{Your turn}
\begin{itemize}
\raggedright
  \item \emph{Say it:} \emph{chīz}
  \item \emph{Say it:} \emph{chīz}, once more
  \item \emph{Say it:} \emph{eh chīz bahut mahiṅgī hai.} — This item is very expensive.
  \item \emph{Recall:} say the Punjabi for to choose, then the Punjabi for to test, then say \emph{chīz} again
  \item \emph{Write it:} copy \textbf{\pa{ਚੀਜ਼}} once, checking the whole word
\end{itemize}

\begin{tcolorbox}[breakable,colback=teal!4,colframe=teal!35!black,title={Before you move on}]
What does \pa{ਚੀਜ਼} mean? (\textquotedblleft{}A thing; an item\textquotedblright{}.) Say it once more.
\end{tcolorbox}
\section[\texorpdfstring{\pa{ਸਾਮਾਨ} (sāmān)}{ਸਾਮਾਨ (sāmān)}]{\pa{ਸਾਮਾਨ} (sāmān) — goods; luggage}

\label{lesson:PA-C277-saman}



\begin{quote}
Before the new one: say the Punjabi for to test, then the Punjabi for a thing.
\end{quote}

\subsection*{You'll want to know: \pa{ਸਾਮਾਨ}}
\textbf{\pa{ਸਾਮਾਨ}} — \emph{sāmān} — \textquotedblleft{}goods; luggage\textquotedblright{}.

\begin{quote}
\pa{ਮੇਰਾ} \pa{ਸਾਮਾਨ} \pa{ਬੱਸ} \pa{ਵਿੱਚ} \pa{ਹੈ।}
\end{quote}

\begin{quote}
\emph{merā sāmān bass vich hai.}
\end{quote}

\begin{quote}
My luggage is on the bus.
\end{quote}

A masculine noun.

\begin{cousinweb}
\textbf{\pa{ਸਾਮਾਨ}} (\emph{sāmān}) is a loan, usually traced to Persian \emph{sāmān}, \textquotedblleft{}goods, things\textquotedblright{}.
\end{cousinweb}

\begin{grammarlens}[title={What you've built}]
One more word for this chapter.
\end{grammarlens}

\subsection*{Your turn}
\begin{itemize}
\raggedright
  \item \emph{Say it:} \emph{sāmān}
  \item \emph{Say it:} \emph{sāmān}, once more
  \item \emph{Say it:} \emph{merā sāmān bass vich hai.} — My luggage is on the bus.
  \item \emph{Recall:} say the Punjabi for to test, then the Punjabi for a thing, then say \emph{sāmān} again
  \item \emph{Write it:} copy \textbf{\pa{ਸਾਮਾਨ}} once, checking the whole word
\end{itemize}

\begin{tcolorbox}[breakable,colback=teal!4,colframe=teal!35!black,title={Before you move on}]
What does \pa{ਸਾਮਾਨ} mean? (\textquotedblleft{}Goods; luggage\textquotedblright{}.) Say it once more.
\end{tcolorbox}
